\documentclass[10pt,a4paper,sans]{moderncv}

%% Hack for "File `l3regex.sty' not found."
\usepackage{expl3}
\expandafter\def\csname ver@l3regex.sty\endcsname{} 

%% ModernCV themes
\moderncvstyle{classic}
\moderncvcolor{blue}
\renewcommand{\familydefault}{\sfdefault}
\nopagenumbers{}

%% Character encoding
\usepackage[utf8]{inputenc}

%% Adjust the page margins
\usepackage[scale=0.75]{geometry}

% For page numerotation lastpage
\usepackage{lastpage}

%% Personal data
\firstname{Bendicht}
\familyname{Grossniklaus}
\title{Microengineer}
\address{Rebhalde 11a}{CH-2555 Brügg}
\mobile{+41 78 654 14 84}
\email{bendicht.grossniklaus@naxys.ch}
%\homepage{www.johndoe.com}
\extrainfo{Born the 17 May 1989\\ Swiss Nationality}
%\extrainfo{Date of birth: 17.05.1989 \\Swiss Nationality\\ Single}
%\photo[64pt][0.4pt]{camipro.jpg}
%\quote{Some quote (optional)}

%%------------------------------------------------------------------------------
%% Content
%%------------------------------------------------------------------------------
\begin{document}
\setcounter{page}{1}
\pagenumbering{arabic}
\rfoot{Bendicht Grossniklaus \hspace{11cm} \thepage\ / \pageref{LastPage}}
%\footnote{Bendicht Grossniklaus}
\makecvtitle

\section{OBJECTIVE}
%\cvitem{}{}%Looking for a job


\section{WORK EXPERIENCE}
\cventry{2021 - 2023}{Electronic development Engineer}{Naxys GmbH}{Prés de l'Envers 2, 2607 Cortébert}{}
{
%\vspace{.3cm}
As consultant, I was working at Sontex SA on the Development of a firmware platform for LoRaWAN radio communication. The goal of this platform was to allow to easily add the LoRaWAN radio function to an existing Sontex-product needing a minimum of changes in the product firmware. Therefore most functionality were implemented on a separate microcontroller and implemented in a way to fit the different Sontex-products. The work consisted of writing specifications, implementing firmware, writing test cases and executing tests for validation.
\vspace{.3cm}
}

\cventry{2020-2021}{Collaborator Kayak and Canoe rental}{GLOBEPADDLER Bielersee AG}{Schlossstrasse 24, CH-2560 Nidau}{}
{
%\vspace{.3cm}
In the summer season 2020 , I was active as collaborator in the Kayak rental GLOBEPADDLER Bielersee AG. The main activity fields were: advice customers for boat type or trip possibilities, teaching kayak courses, transport kayak-canoe and rented material with the company's car, minor reparations on the rented material and the rental station.
\vspace{.3cm}
}


\cventry{2019 - 2020}{Electronic development Engineer}{Société industrielle de Sonceboz SA}{Rue Rosselet-Challandes 5, 2605 Sonceboz-Sombeval}{}
{
%\vspace{.1cm}
In the \textit{Mechatronics} division of the SONCEBOZ GROUP, I worked in a team to implement the software for a new family of actuators for hydraulic valves. As the exclusive electronic engineer of the project team, I was in charge of the software and electronics of an actuator driving a Continously Variable Transmission (CVT) of agriculture engines. In close collaboration with the customer and the project manager, we enhaced the CVT actuator from a prototype to a pre-series status within this year.
\vspace{.3cm}
}

\cventry{2016 - 2019}{Electronic development Engineer}{Sonceboz Microtechnique Boncourt SA}{Rue des Boulaies 2, CH-2926 Boncourt}{}
{
%\vspace{.1cm}
Following my Master Thesis in Munich, I was mandated as electronic development engineer in the \textit{Truck and Offroad} division of the SONCEBOZ GROUP. Thanks to my knowledge of the CPM products and the good relationship with the colleagues in Munich, Sonceboz and Boncourt, I participated as a linking part to bring the CPM products to maturation for series production and gave support to potential customers. During this time I did software development such as implementation of position control, introduction of an analog speed throttle and resolutions of issues. Further, I introduced new PCB versions for the CPM products, being in charge of organizing and supervising prototype-mounting with the Manufacturing Unit. Finally, I gave technical support to customers around the globe and to colleagues on the production-line (investigating of product and production issues on the line).
\vspace{.3cm}
}

\cventry{2015-2016}{Student Internship}{Société industrielle de Sonceboz SA}{Rue Rosselet-Challandes 5, 2605 Sonceboz-Sombeval}{}
{
%\vspace{.1cm}
For my end-of-studies internship, I was mandated by the \textit{Truck and Offroad} division of the SONCEBOZ GROUP to investigate the feasibility of measuring the temperature of a BLDC motor using the variation of its coil resistance. Using the applied voltage and the current measurement, the resistance variation of the motor phases could be computed and thereby, the temperature variation estimated. Implementing the software and leading measurement campaigns, I could show that the method allowed a rough estimation of the coil temperatures. At the end of this six-months internship, I wrote a laboratory journal in perspective to a potential patenting of the method.
\vspace{.3cm}
}

\cventry{2014}{Civilian Service}{Ecological Engineering Laboratory (ECOL)}{EPFL, CH-1015 Lausanne}{}
{
%\vspace{.1cm}
Working in a team of civilists, we improved the robustness and added new functionalities to an autonomous motorized catamaran based on a NI-Computer programmed in LabVIEW. Further, we helped to collect data for the doctoral candidates on the lake Léman using the catamaran, during the measuring campaigns.
\vspace{.3cm}
}

\cventry{2012 - 2013}{Freelance translator}{SABAG LUZERN AG}{CH-6023 Rothenburg}{}
{
%\vspace{.1cm}
As a freelance translator, I did a German to French translation of a tutorial introducing a new logistic program to the employees.
\vspace{.3cm}
}

\cventry{2008 - 2009}{Civilian Service}{CH-4000 Basel Hörnli, CH-2608 Prés-de-Cortébert, CH-2503 Biel/Bienne}{}{}
{
%\vspace{.1cm}
Different mandates with a total duration of seven months in the following fields : landscape gardening, work in an organic farm, building dry stone walls
\vspace{.3cm}
}


\section{EDUCATION}
\cventry{2014 - 2016}{Master}{Ecole Polytechnique Fédérale de Lausanne (EPFL)}{CH-1015 Lausanne}{}{Micro-engineering, specialization in Robotics and Autonomous Systems (RSA)}  % arguments 3 to 6 can be left empty
\cventry{2009 - 2014}{Bachelor}{Ecole Polytechnique Fédérale de Lausanne (EPFL)}{CH-1015 Lausanne}{}{Micro-engineering}  % arguments 3 to 6 can be left empty
\cventry{2008}{High school diploma}{Gymnasium Alpenstrasse}{CH-2502 Biel-Bienne}{}{Physics and applied mathematics, bilingual (French-German)}  % arguments 3 to 6 can be left empty


\subsection{PROJECTS}
%\subsection{Master Diploma Project}
\cventry{2016}{Master Thesis: Design of an external BLDC motor with integrated electronics for 48V on-board power supply}{Sonceboz SA - Integrated Actuators Laboratory (LAI)}{CPM Compact Power Motion GmbH, Feringastraße 11A, 85774 Unterföhring, Deutschland}{}
{
	%\vspace{.05cm}
	As a collaboration between the SONCEBOZ GROUP and LAI-EPFL, my Master Thesis took place in Munich, where I was immersed in the Spin-Off Compact Power Motion GmbH (CPM). The goal of the work was to design a motor of the same type as the existing products of CPM (BLDC motors with external rotor), but with a two-times higher power and torque. The main topics of the work were the design of the motor geometry and its optimization using 2D FEM simulations (Flux2D), the evaluation of the thermal performance using FEM simulations (SolidWorks) and the implementation of the chosen solution as a model for the Software-In-The-Loop Tool (Matlab/Simulink). The outcome was a solution that fulfilled the defined requirements considering torque, power and control of the motor (driving the model with the SIL). Nevertheless the thermal behavior of the newly-designed motor was worse than the existing CPM products due to the bigger diameter of the inner stator. During the six month of internship, I met many people of the SONCEBOZ GROUP on the three sites Sonceboz (R+D), Boncourt (Production) and Munich (R+D) to get the information I needed to accomplish my Master Thesis.
	\vspace{.3cm}
}

%\subsection{Semester projects}
\cventry{2015}{Semester Project II: Automation of a laboratory coil winding machine}{Integrated Actuators Laboratory (LAI)}{Rue de la Maladière 71 B, CH-2002 Neuchâtel}{}
{
	%\vspace{.1cm}
	The goal of my second semester project was to automatize a semi-automatic coil winding machine used by the LAI for prototyping. I did this using a commercial Arduino-board and a linear stepper motor. The stepper motor is synchronized with the rotary motor using it's encoder signals and an LCD with five push-buttons for user entries.
	\vspace{.3cm}
}

\cventry{2014-2015}{Semester Project I: Development of a soft jellyfish type underwater robot}{Laboratory of Intelligent Systems (LIS)}{EPFL, CH-1015 Lausanne}{}
{
	%\vspace{.1cm}
	The goal of my first semester project was to design a jellyfish-like robot using Dielectric Elastomer Actuators (DEAs). In order to find the best design, I did two models in MATLAB, one for the propulsion force and one for the DEA deformation. My design was built in a specialized laboratory by the doctoral student supervising my work. I mad the test on the final prototype measuring the force, the deflection and the power consumption of the actuated jellyfish in water. With this work, we could show that with its high energy density and its fast actuation response, DEA is a promising technology for soft robotics.
	\vspace{.3cm}
}

\section{TECHNICAL SKILLS}
\cvitem{Programming Languages}
{
Good knowledge : C with MISRA-C Standard, VECTOR CAPL, MATLAB/Octave scripts
Basic knowledge : ASSEMBLER, C++
}

\smallskip

\cvitem{Software Tools}
{
CodeWarrior/Code Composer Studio (Eclipse IDE), SpiraTeam, Altium Designer, Vector CANalyzer, SolidWorks (with PDM), MATLAB/SIMULINK, Flux2D/3D and LabVIEW.
}

\smallskip

\cvitem{Handcraft}
{
Solder electronic circuits for prototyping, building small test or measurement setups and traditional manufacture know-how (drilling and milling).
}




\section{LANGUAGES}
\cvitem{German}{Bilingual (Swiss German Mother Tongue)}
\cvitem{French}{Bilingual (Obligatory school)}
\cvitem{English}{Good knowledge (Level B2)}
\cvitem{Spanish}{Good oral understanding}




\section{HOBBIES AND INTERESTS}
\cvitem{Sports}{Swiss junior national team kayak slalom (2003-2007), J+S guide in kayak since 2007, leading a regular training to club youth since 2017 capoeira, kitesurfing, snowboarding and biking}

\smallskip

\cvitem{Handicraft work}{Transformation of a van into a camper, building an exploiting an aquaponic-sytem, permaculture and traditional gardening}
%\vspace{0.1cm}
%\cvitem{\textbf{Robopoly}}{Member of the EPFL robotic association, participation in the annual robot contest of the school (2011, 2012, 2014).}

\bibliography{publications}
\end{document}